% Include-only section; requires amsmath, amssymb, amsthm, booktabs.
\section{A productive-adjacency obstruction for rotating pop}
\label{sec:rotating-obstruction}

The archived packet \texttt{MP0\_Transfer\_215.zip} contains a physically
validated experimental 215-state table, \texttt{research/rotate-pop215/R.native.tm},
whose refutational completeness is explicitly left open. Its command
\[
 \mathsf R(t,[x;v],d,h)=(x,v,d,t)
\]
halts instead when $h=1$ or $x=1$. Here $[x;v]=P(x,v)$, with Cantor
pairing $P$, and $\operatorname{unpair}(0)=(0,0)$. The following argument
shows that the missing compiler cannot be repaired while keeping this
interface and the literal AX row order. It uses no consistency assumption,
no model of set theory, and no theoremhood invariant.

Put $I(a,b)=4P(a,b)$, $Q_a b=4P(a,b)+2$, $E(a,b)=4P(a,b)+3$,
$M(a,b)=4P(a,b)+1$, $F=1$, and $NX=I(X,F)$. Define
$\exists_aX=NQ_aNX$. The actual AX output is
\[
 (G,[r_1,r_2,r_3,r_4,r_5,r_6,r_7,L,Q_Uh,t;w],0,0),
\]
where $U=d$, $A=d+x$, $B=h+y$, $C=c$, the $r_j$ are exactly
\begin{align*}
 r_1&=\exists_AE(A,B),\\
 r_2&=I(Q_A I(B,C),I(Q_AB,Q_AC)),\\
 r_3&=I(E(A,B),I(M(C,A),M(C,B))),\\
 r_4&=I(Q_AQ_BC,Q_BQ_AC),\\
 r_5&=I(E(C,A),I(M(C,B),M(A,B))),\\
 r_6&=I(\exists_AQ_AC,C),\\
 r_7&=I(E(C,A),I(E(C,B),Q_{M(A,B)}E(A,B))),\\
 L&=I(I(I(C,U),A),I(I(A,C),I(B,C))).
\end{align*}
The gate $G$ is either zero or the literal full bundle
\[
 G=I(I(V,J),F),\qquad V=I(M(A,B),Q_CM(A,B)).
\]
The contents of $J$ are immaterial to the argument: the full inherited
term is retained. Thus the result does not weaken the underlying theory.

Call a pair $(m,a)$ \emph{productive} when
\[
             \exists b>0\quad m=I(a,b).
\]
For a command store, consider the ordered spine
$h,t,x_1,x_2,\ldots$, where $x_i$ are successive heads of its encoded
stack. A spine is \emph{blocked} when none of its entries is $1$ and
none of its adjacent pairs is productive. The all-zero spine is blocked.
This definition is numerical; it does not require that an entry decode
to a conventional formula.

\paragraph{The AX adjacency lemma.}
Every newly introduced adjacent pair of the AX spine
\[
 0,G,r_1,r_2,r_3,r_4,r_5,r_6,r_7,L,Q_Uh,t,x_1,\ldots
\]
is unproductive. The existing suffix beginning with $t$ is unchanged.
For the pair beginning with zero this follows from the exact implication
alias $I(a,b)=0\Longleftrightarrow a=b=0$. For the remaining pairs,
injectivity of $I$ reduces a possible match $m=I(a,b)$ to equality
between the antecedent of $m$ and $a$. The following table gives the
contradictory outer tags; nested comparisons use injectivity once more.

\begin{center}\small
\begin{tabular}{lll}
\toprule
Pair & Required equality & Contradiction modulo $4$\\
\midrule
$(G,r_1)$, enabled gate & $V=Q_A N E(A,B)$ & $0\ne2$\\
$(r_1,r_2)$ & $Q_A N E(A,B)=r_2$ & $2\ne0$\\
$(r_2,r_3)$ & $Q_A I(B,C)=r_3$ & $2\ne0$\\
$(r_3,r_4)$ & $E(A,B)=r_4$ & $3\ne0$\\
$(r_4,r_5)$ & $Q_AQ_BC=r_5$ & $2\ne0$\\
$(r_5,r_6)$ & $E(C,A)=r_6$ & $3\ne0$\\
$(r_6,r_7)$ & $Q_A N Q_AC=E(C,A)$ & $2\ne3$\\
$(r_7,L)$ & $E(C,A)=L$ & $3\ne0$\\
$(L,Q_Uh)$ & $I(I(C,U),A)=Q_Uh$ & $0\ne2$\\
$(Q_Uh,t)$ & $Q_Uh=I(t,b)$ & $2\ne0$\\
\bottomrule
\end{tabular}
\end{center}
For a disabled gate $G=0$, the zero-pair argument applies. All generated
entries are different from $1$: the gate and eight rows have tag zero,
and $Q_Uh$ has tag two. This includes the possible alias $L=0$; no
positivity of $L$ is assumed. The inequality $U\le A$ is not needed.

\paragraph{Preservation theorem.}
Every finite word of the rotating-pop commands, MP and AX preserves a
blocked spine and never halts when begun at a blocked spine. In
particular, this holds for every finite word from $t=h=w=0$, with
arbitrary additive counter and arbitrary raw command parameters.

\emph{Proof.} AX preserves the property by the adjacency lemma.
Rotating pop maps the spine $h,t,x_1,x_2,\ldots$ to
$t,x_1,x_2,\ldots$, a suffix. Its two tested entries are not $1$.
For MP put $y=d+x$. If it succeeds, $h=I(t,y)$, so blockedness forces
$y=0$. The returned spine is $0,0,0,x_1,x_2,\ldots$. If it fails,
the returned spine is $0,0,x_1,x_2,\ldots$. Every new pair beginning
with zero is unproductive, and every new entry is zero. In either
case blockedness is preserved. Induction on the word proves the claim.
This proof also covers every numerical zero-stack alias.\hfill$\square$

\paragraph{The actual dirty-tail initialization.}
The preceding argument has a finite-prefix form. For any fixed raw tail
$J$, maintain a finite explicit head list $X$ over $J$ such that the
finite spine $[h,t]+X$ is blocked. AX adds ten explicit heads, R consumes
one, and MP adds either zero heads or one zero head. Whenever the
current explicit head list is nonempty, the same preservation proof
applies. Thus an external scheduler credit bound protects \emph{blocked}
heads just as it protects theorem-valued heads; the latter interpretation
is unnecessary here.

For the archived R table, the supplied exact blank startup reaches the
interior of the unchanged AX at source ordinal $55$, with registers
$[1,11,2,0,0,0,0,0,0]$. Its recorded AX-prefix match completes to next
candidate $1$, empty proof pointer, $t=d=h=0$, and
\[
 [r_1(0,0,0),\ldots,r_7(0,0,0),L(0,0,0,0),Q_0 0,0;4].
\]
The displayed ten-head prefix is blocked by the lemma, independently
of the tail $4=P(1,1)$. The unchanged wire arities are one field for
R (tag zero), two fields for MP (tag one), and three fields for AX
(tags at least two). The inherited all-natural Cantor-diagonal credit
proof assigns weights $-1,0,10$ and proves that every schedule prefix
from candidate $1$ has weight at least $-2$. Consequently at least
eight explicit blocked heads remain throughout the actual execution.
The buried $1$ heads are never reached. The source therefore never
halts, by an ordinary proof formalizable in ZFC. Combining the supplied
positive-duration primitive simulation and the complete blank prefix
gives the same unconditional nonhalting statement for the exact R table.

This last exact-table conclusion is conditional on those independently
specified physical certificates; the new proof here establishes the
command-level obstruction and its protected-prefix transport. It does
not rerun the large physical certificates. The checker accompanying this
section verifies all ten literal adjacency obstructions and the exact
numerical gate/row formulas used in this proof. It does not claim to
be a proof-assistant kernel check.

\paragraph{A fixed theorem defeats the proposed strong compiler.}
The calibrated numerical calculus proves
\[
                   X=I(0,Q_0 0)=I(0,2)=20.
\]
Indeed truth gives zero, GEN gives $Q_0 0=2$, and the derived
propositional K schema gives $2\to(0\to2)$; one MP gives $X$.
The dominated calculus inherits that K schema from its established
propositional recovery certificate. A strong block for $X$, applied
to the all-zero store, would have to halt or return
$(20,0,0,0)$. The preservation theorem excludes halting. The alleged
return has spine $0,20,0,0,\ldots$, containing the productive adjacent
pair $(20,0)$ because $20=I(0,2)$ and $2>0$. Hence the return is
also impossible. No strong compiler for all numerical theorems exists,
even for this fixed elementary theorem; this conclusion does not use
any consistency or incompleteness assumption.

\paragraph{Consequence for refutational completeness.}
More generally, no block can take a blocked all-zero store to a
return state with top $1$ or to a halt. If a
refutation-complete compiler for the calibrated ZFC calculus existed
and were provable in ZFC, the unconditional nonhalting proof would
yield $\operatorname{Con}(\mathrm{ZFC})$ in ZFC. Under the usual
consistency assumption this contradicts the second incompleteness
theorem. The obstruction is specifically caused by preserving the
literal AX order while replacing transfer-pop with rotating pop.
Changing that order, emitting an additional bridge pair, or changing
the return map would need a new analysis.
